% !TEX root = _main.tex
%1
\section{はじめに}

近年インターネットの普及に伴い音楽の楽しみ方が大きく変化している．
従来ではCDを購入してプレイヤーで楽しむことが一般的であったが，近年はインターネット配信で楽しむ人が増加している．
各音楽配信サイトなどでは，ユーザの嗜好に合った楽曲の検索を容易にするために楽曲を音楽ジャンル別で分類されることが多くみられる．

音楽ジャンル解析は古くから数多くの研究が行われており，近年では主要な研究テーマの１つとして扱われている．
しかし，楽曲のジャンルは時系列で変わっていくのではないだろうか？
%
\red{例えば，}Mrs. GREEN APPLEの「ライラック」\cite{lilac}%があげられる．
%この楽曲
\red{で}は，冒頭約10秒間でギターにエフェクトを加えた電子音が流れる．
しかし，その後同じフレーズをドラムやギターの音色で構成する．
これにより，電子音のような音色では感じることのなかったロックミュージックの印象を抱かせるものとなっている．
このように，聴取する箇所によって抱く印象が異なる楽曲が存在すると考える．

%しかし，これらの大半は１つの楽曲を1つの音楽ジャンルに分類するものである．
%このような音楽ジャンルの分類方法では，楽曲内における音楽ジャンルの変化の識別が困難であると考える．

そこで本研究では，数ある音響特徴量の中でも\red{音色に着目し，対象楽曲の楽器構成から\footnotemark}音楽ジャンルの変化を時系列で表示するシステムを提案することを目的とする．
\footnotetext{\red{「音色」に着目して（いろいろ検討した結果，）楽器構成を手掛かりにジャンル分析をするのがよいとした，その思考プロセスを２章でしっかり語るべし．}}



%2
\section{先行研究}
現在まで，音楽ジャンル解析に関する研究は数多く行われている．楽曲のパート編成やピッチ，音長，リズムパターン，音色，リズム特性などといった特徴量を対象とした研究がある\red{（［文献URL追加］）}．
\red{ジャンル推定をするのに，なぜ楽器編成に着目するのか？従来研究の知見とともにりささんの主張をここで言う．}


\subsection{\red{楽器構成に着目したジャンル推定}}


土橋らは音響信号の中でもベースパートに着目し，ベースパートの特徴量，音色，リズムを用いた音楽ジャンル推定を行っている\cite{bassline}．また，鈴木らは旋律パートのみを用いて音楽ジャンル判別を行うための新たな特徴量を提案し，その有効性を検証している\cite{melody}．

Basiliら\cite{basili}はMIDIデータから楽曲のパート編成やピッチ，音長といった高次の特徴量を抽出し，機械学習アルゴリズムを用いて音楽ジャンルを分類している．

また，楽器構成の抽出を応用した研究も行われている．

卓見ら\cite{takumi}らは，楽器編成を自動認識し，それを用いて楽曲を探索するための手法を開発している．
楽器編成の認識方法には畳み込みニューラルネットワーク（CNN）を使用し，既存の方法よりも高精度な認識が達成された．

Yoonchang Hanら\cite{han}は，畳み込みニューラルネットワーク（ConvNet）を使用して，ポリフォニック音楽の音源から主要楽器を識別する方法を提案している．
なお，ポリフォニック音楽とは複数の楽器が同時に演奏される音楽である．
畳み込み層とプーリング層を繰り返し使用したConvNetを採用し，結果として，ConvNetは精度と再現率において従来の手法を上回っている．


\subsection{\red{時系列の推移に対応した類似演奏の探索}}
\red{
	query-by-conducting，mazurka projectみたいなことをしたい．という話がごっそりぬけている．
	図を引用してきていいので，自分が「時系列に対応したジャンル推定」の出力がどういうイメージを持っているのかを示す．
	音楽の類似度（music similarity）をどうやって計算するか，という話も本当はしたい．
	少なくとも質問されたら答えられるように！
}

\subsection{\red{本研究の狙い}}
\red{
	本研究が目指す世界（達成したいこと）が何かを明確にする．
	ポイントは，(1)楽器編成に注目するとジャンル分析がうまくいく（という期待があること）
	(2) 時系列で区切って，区間別にジャンルっぽさの変化の様子を可視化したい
}

上記で記した研究は１つの楽曲に対して１つのジャンルを判別するものであり，従来はこの判別方法が典型であった．しかし，本研究では楽曲を時間で区切ることにより，楽曲内の音楽ジャンルの変化を可視化することが可能となる．
